\subsection{A finite counterexample with a certified gain}

The counterexample can use just the $224$ odd primes $q\le Y=1427$.
Let $B$ denote these primes, $Q=\prod_{q\in B}q$, $S=\{2\}\cup B$, and
$C$ the Rankin subset of the $S$-smooth integers. Write $a=v_2(n)$ and
$b_q=v_q(n)$. For an old point $n\in C$, delete it under either of the
following conditions.
\begin{itemize}
\item All its odd exponents $b_q$ equal $1$ and $a>0$.
\item Exactly one odd exponent $b_q$ differs from $1$, and either $b_q=4$,
or $b_q=3$ and $a\in2K$, where $K=2+3\{0,1,3,4\}=\{2,5,11,14\}$ and $2K=\{2k:k\in K\}$.
\end{itemize}
Insert the points $2^kQ$, $2^kQq$ and $2^kQq^3$ for every $q\in B$ and $k\in K$.
All four markers are $2\bmod 3$, so these points are disjoint from $C$.

\paragraph{The local obstruction.}
Let $G$ be the old ternary exponent set.
If $u,v\in G$ and $u,2,v$ is an arithmetic progression, then
$\{u,v\}=\{0,4\}$ or $\{1,3\}$. If $u,v,2$ is a progression, then
$u=2v-2$ and $v\in\{1,3\}$. To verify the last assertion, write $v$
in ternary. If the last digit is $1$, subtracting $2$ from $2v$ leaves
each higher digit doubled, so all higher digits must vanish. If the
last digit is $0$, let $j\ge1$ be the position of the first nonzero
digit. Borrowing to subtract $2$ would create a digit $2$ at position
$1$ if $j>1$. Thus $j=1$ and all higher digits must vanish, again by
doubling. The two possible $v$ do work.

Consider a progression with only the insertion $2^kQq$ among its three
members. All its other odd exponents are $1$, by progression avoidance
in $G$. If the insertion is an endpoint, the middle $q$-exponent is
$1$ or $3$ by the local obstruction. In the first case, the middle
point has all odd exponents $1$ and positive $2$-exponent, since it
averages the positive $k$ and a nonnegative exponent. In the second
case, the other endpoint has $q$-exponent $4$. Both cases give a
deleted point. If the insertion is the midpoint, the outside
$q$-exponents are $0,4$ or $1,3$. The point with $q$-exponent $4$
was deleted. In the other pair the point with $q$-exponent $1$ survives
only at $a=0$, which forces the point with $q$-exponent $3$ to have
$a=2k\in2K$, also deleted. For a single insertion $2^kQ$, both
old points must have all odd exponents $1$, and can survive only at
$a=0$. The values $0,0,k$ cannot form a progression. For a single insertion $2^kQq^3$, the exceptional odd exponent is $4\in G$. Every odd coordinate of all three points then belongs to $G$ and is constant. Both old points would have exactly the deleted odd pattern with exceptional exponent $4$.

If two points were inserted, their $2$-exponents are $2\bmod3$, forcing
the third exponent to have that residue in any of the three positions.
No member of $G$ has this residue. With three insertions, every odd exponent
is in the progression-free set $\{1,2,4\}$ and the $2$-exponent is in the progression-free set $K$ (an affine image of a subset of $G$). Each coordinate
must therefore be constant in a progression, so all terms coincide.
This proves avoidance for every rational ratio.

\paragraph{Size of the improvement.}
Let $S_2=\sum_{a\in G}2^{-a}$. Exponents below $27$ use only the
three lowest ternary digits. Bounding the remaining terms by
$\sum_{a\ge27}2^{-a}=2^{-26}$ gives
\[
 S_2\le U=(1+2^{-1})(1+2^{-3})(1+2^{-9})+2^{-26}.
\]
Write $u=\sum_{k\in K}2^{-k}$ and $v=\sum_{k\in K}2^{-2k}$.
Let $\Delta W$ denote the change in reciprocal mass on smooth numbers.
In units of $1/Q$, the insertions contribute $u(1+\sum_{q\in B}(q^{-1}+q^{-3}))$,
the common deletion costs $S_2-1$, and the two kinds of
single-coordinate deletions cost at most $v/q^2$ and exactly $S_2/q^3$,
respectively. The first estimate also counts any $2k$ absent from $G$. Consequently
\[
 Q\Delta W\ \ge\ u-(U-1)
     +\sum_{q\in B}\left(\frac{u}{q}-\frac{v}{q^2}-\frac{U-u}{q^3}\right)
       >\frac9{100000}.
\]
The last strict inequality is checked in exact rational arithmetic by
\texttt{check\_rankin\_finite.py}. The script sieves the primes, verifies
$Q<3\cdot10^{590}$, and checks
\[
          \frac{Y}{Y+1}\prod_{p\le Y}(1-p^{-1})>\frac1{15}.
\]
The outside-prime density factors are at least $1-p^{-2}$, so
their product is at least $\prod_{n>Y}(1-n^{-2})=Y/(Y+1)$.
The density $c_H$ of the permitted rough parts is therefore greater
than $1/15$. Multiplying the smooth gain by $c_H$ yields the explicit
bound for the natural density of the full construction,
\[
              d(A)-d(R) > \frac{9}{100000\cdot15\cdot3\cdot10^{590}}
                           > 10^{-596}.
\]
